\documentclass[11pt]{article}

\usepackage[margin=1in]{geometry}
\usepackage{booktabs}
\usepackage{graphicx}
\usepackage{amsmath}

\setlength{\emergencystretch}{2em}

\title{Weekly Report: Week 3}
\author{}
\date{2026-08-12}

\begin{document}
\maketitle

\section{Scope}

This report covers the pair-adaptive attribute-selection stage (Step~1.5) of the
LLM-assisted entity resolution pipeline, measured on Fodors--Zagat and DBLP--ACM: the
method, the experimental design the comparisons require, the results, and what those
results do and do not establish.

All numbers come from runs recorded on 2026-08-12 and 2026-08-18. The matcher LLM is
\texttt{deepseek-\allowbreak v4-\allowbreak flash-\allowbreak 0731} (served through
OpenRouter) and the embedding model is
\texttt{bge-\allowbreak large-\allowbreak en-\allowbreak v1.5}. Results predating this
report are not comparable and are not quoted: they were produced under a different matcher
LLM and before several defects in the experimental setup were repaired. Amazon--Walmart is
covered only by its no-selection baseline (Section~\ref{sec:amazon}); the remaining
configurations were still running when this report was written.

The headline result is that on DBLP--ACM, showing the matcher a small, pair-dependent
subset of attributes beats both showing it everything and showing it a fixed subset chosen
with labelled data --- but only under a selection rule that can express small subsets, and
the rule matters as much as the method does.

\section{Pipeline Overview}

The pipeline runs in five stages, each timed and recorded in a per-run JSON summary under
\texttt{logs/<dataset>/runs/}.

\textbf{Load and normalize.} Both tables and the ground truth are read, ID columns are
dropped, and every remaining field is normalized: HTML entity references are decoded, text
is Unicode-decomposed with combining marks stripped, and address and directional
abbreviations are standardized. This matters more than it appears: \texttt{ACM.csv} stores
461 HTML entity references (\texttt{\&\#246;}, \texttt{\&mdash;}) where \texttt{DBLP.csv}
stores none, so without decoding, the same author name is two unrelated strings across the
two sources.

\textbf{Embed and block.} Each attribute column is encoded independently with SentenceBERT,
the column vectors are concatenated and L2-normalized into one embedding per record, and
multi-probe locality-sensitive hashing retrieves the top-$k$ most similar Table-B candidates
for every Table-A record. On DBLP--ACM this reduces 6{,}001{,}104 possible pairs to 13{,}080
candidates while retaining all 2{,}224 ground-truth matches.

\textbf{Step 1 --- global tuple condensation (optional).} One attribute subset is chosen
once for the whole dataset and applied to every record. The tables may optionally be
re-blocked on the condensed attributes; see Section~\ref{sec:reblock}.

\textbf{Step 1.5 --- pair-adaptive attribute selection (optional).} A small sample of
candidate pairs is profiled by the LLM for per-attribute importance, and a hybrid
predictor/retriever model transfers a \emph{different} attribute subset to every candidate
pair. Blocking is unchanged, so the candidate set is identical to the no-selection baseline.

\textbf{Match.} One LLM call decides Yes/No per candidate pair, using only the attributes
selected for that pair.

\textbf{Metrics and logging.} Precision, recall, and F1 are computed end-to-end, against all
ground-truth matches, so pairs lost at blocking count against recall.

\section{Methodology}
\label{sec:methodology}

\subsection{Step 1.5: Pair-Adaptive Attribute Selection}

Step~1.5 has a profiling stage and a transfer stage.

\textbf{Profiling.} A label-free sampler selects a small set of candidate pairs (20 on
Fodors--Zagat, 50 on DBLP--ACM) and the LLM is asked, for each, which attributes matter for
deciding whether the two records describe the same entity. Ground truth is not consulted at
any point. The \texttt{diversity} sampler is used throughout this report, holding the
sampler fixed so that other variables can be read.

\textbf{Transfer.} For every candidate pair, six cheap similarity features are computed per
attribute --- token overlap, edit similarity, exact match, semantic similarity, numeric
similarity, and a numeric-type indicator --- without calling the LLM. Two models then
predict per-attribute importance from those features: a Ridge regressor fit on the profiled
pairs, and a nearest-neighbour retriever that averages the importance vectors of the five
most feature-similar profiled pairs. Their outputs are blended equally and normalized to sum
to one, and a selection rule decides how many of the ranked attributes to keep.

The profiling stage costs about 1\% of a run's tokens. Its purpose is to buy a rule that can
be applied to the remaining 99\% without further LLM calls.

\subsection{Two Ways to Elicit Importance}
\label{sec:scoring}

The profiling prompt is one of the two variables this report isolates. Two formulations were
compared, identical in framing and guardrails and differing only in the answer requested.

\textbf{Ordinal.} The LLM assigns each attribute an integer score from 0 (not useful) to 3
(decisive), scoring each attribute independently.

\textbf{Decisive.} The LLM returns the smallest subset of attributes sufficient to make the
decision, weighing attributes against one another in a single judgement.

The motivation for the second formulation was that the recorded 0--3 scores were not being
used discriminatively. Aggregated over all runs, DBLP--ACM used score 0 on 1\% of
judgements, collapsing a four-point scale to three; Fodors--Zagat used score 3 on 55\%. An
importance vector that is nearly flat cannot express that one attribute matters more than
another, whatever selection rule consumes it.

\subsection{How Many Attributes to Keep}
\label{sec:policy}

Profiling and transfer produce a ranked importance vector per pair. A separate rule turns
that ranking into a subset, and it is the second variable isolated here.

\textbf{Cumulative} keeps attributes in descending importance until their shares sum past a
fixed target, by default 0.8. It is the pipeline's historical default and is not scale-free:
importance is normalized to sum to one, so on a wide schema each attribute's share shrinks
and a fixed target becomes unreachable. On Amazon--Walmart's 13 attributes the 0.8 target
needs 7.7 attributes on average against a cap of 5, so the cap binds for every pair and the
rule cannot express ``this pair needs fewer fields'' at all.

\textbf{Gap} cuts at the largest drop in the ranked importance profile. It has no threshold
to tune and does not saturate on any of the three schemas.

Both rules are subject to a floor and a cap on the number of attributes kept. The floor was
2 by default and is set to 1 wherever a run is allowed to reduce a pair to a single
attribute.

\section{Experimental Design}
\label{sec:design}

\subsection{Why Comparisons Are Paired}
\label{sec:paired}

The matcher LLM is not deterministic, and on these datasets its variation is large enough to
swamp the effects being measured. Running the no-selection baseline twice on Fodors--Zagat,
with identical prompts and identical attributes, produced F1 of 0.8917 and 0.9596 --- a gap
of 0.068, from 19 of 2{,}665 answers changing, 18 of them from Yes to No. Precision moved
from 0.823 to 0.947 while recall was unchanged.

The same measurement on DBLP--ACM gives a much smaller floor: two independent draws of the
no-selection baseline produced F1 of 0.8706 and 0.8613, a gap of 0.0093. This is expected ---
DBLP--ACM has 2{,}224 ground-truth matches against Fodors--Zagat's 110, so the same absolute
number of flipped answers perturbs the metrics far less.

Comparing two methods by their F1 values from separate runs therefore measures the matcher's
variance more than it measures the methods. Two consequences follow.

First, the matcher response cache is enabled for all comparisons. Where two methods select
the same attributes for a pair, the pair yields one shared answer and cannot contribute a
spurious difference. Section~\ref{sec:mechanism} shows this working: across the 9{,}803 pairs
where two very different methods happened to select the same single attribute, they recorded
identical false-positive counts.

Second, methods are compared by McNemar's test on the pairs where they actually differ,
rather than by comparing aggregate F1. This is a paired test on the same candidate set and
the same matcher draw, and it is the only comparison in this report that carries a
significance claim.

A caveat this design imposes: because repeated runs of one configuration replay the cache,
they are not independent measurements and their standard deviation is meaningless.
Replication across seeds was dropped in favour of the paired design.

\subsection{Holding the Candidate Set Fixed Across Methods}
\label{sec:reblock}

Step~1.5 leaves blocking untouched, so its candidate set is identical to the no-selection
baseline. Step~1 condenses the record and can optionally re-block on the condensed
attributes, which would change which pairs are compared. It is therefore run with
re-blocking disabled (\texttt{--step1-no-reblock}) throughout, so that every configuration is
scored over the same 13{,}080 candidate pairs containing the same 2{,}224 ground-truth
matches, and attribute choice is the only difference between them.

\section{Results}

\subsection{Fodors--Zagat: A Saturated Benchmark}

\begin{table}[h]
\centering
\resizebox{\textwidth}{!}{%
\begin{tabular}{lrrrrrrrr}
\toprule
Method & TP & FP & FN & F1 & Precision & Recall & Attrs/pair & Prompt tokens \\
\midrule
Step 1.5, decisive        & 106 &  4 & 4 & 0.9636 & 0.9636 & 0.9636 & \textbf{2.68} & 652{,}422 \\
Step 1.5, ordinal         & 108 &  6 & 2 & \textbf{0.9643} & 0.9474 & \textbf{0.9818} & 4.00 & 705{,}395 \\
No selection (\texttt{full}) & 107 &  6 & 3 & 0.9596 & 0.9469 & 0.9727 & 5.00 & 748{,}469 \\
Step 1, supervised (10 labels) & 106 &  5 & 4 & 0.9593 & 0.9550 & 0.9636 & 2.00 & 605{,}836 \\
\bottomrule
\end{tabular}}
\caption{Fodors--Zagat, \texttt{class} excluded. 2{,}665 candidate pairs, 110 ground-truth
matches. Metrics are end-to-end. All rows use the cumulative rule at its 0.8 default.}
\end{table}

Every pairwise McNemar test returns $p = 1.00$. The largest disagreement between any two
methods is 9 pairs out of 2{,}665, split almost evenly in both directions.

This is a property of the benchmark, not of the methods. Doing nothing at all already
reaches F1 0.9596 with 107 of 110 matches recovered and 6 false positives. There is no
headroom in which a method can distinguish itself, and with 110 ground-truth pairs the
matcher's own instability moves F1 further than any method does.

What remains measurable is cost. The decisive formulation shows the matcher 2.68 attributes
per pair instead of 5.00, a 46\% reduction, changing 4 of 2{,}665 decisions, two in each
direction. The defensible claim from this dataset is that adaptive selection is free: it does
not cost accuracy, and it does not buy any.

\subsection{DBLP--ACM: Where the Methods Separate}
\label{sec:results-dblp}

\begin{table}[h]
\centering
\resizebox{\textwidth}{!}{%
\begin{tabular}{llrrrrrrr}
\toprule
Method & Rule & Attrs/pair & TP & FP & FN & F1 & Precision & Recall \\
\midrule
Step 1.5, decisive  & gap            & 1.42 & 2197 & \textbf{255} & 27 & \textbf{0.9397} & \textbf{0.8960} & 0.9879 \\
Step 1, supervised (200 labels) & --- & 1.00 & 2198 & 339 & 26 & 0.9233 & 0.8664 & 0.9883 \\
Step 1.5, decisive  & cumulative 0.8 & 2.24 & 2214 & 535 & 10 & 0.8904 & 0.8054 & 0.9955 \\
No selection (\texttt{full}) & ---   & 4.00 & 2216 & 651 &  8 & 0.8706 & 0.7729 & \textbf{0.9964} \\
Step 1.5, ordinal   & cumulative 0.8 & 3.77 & 2215 & 699 &  9 & 0.8622 & 0.7601 & 0.9960 \\
Step 1, supervised (10 labels)  & --- & 2.00 & 2211 & 739 & 13 & 0.8547 & 0.7495 & 0.9942 \\
\bottomrule
\end{tabular}}
\caption{DBLP--ACM. 13{,}080 candidate pairs, 2{,}224 ground-truth matches. Metrics are
end-to-end and every configuration is scored over the same candidate set
(Section~\ref{sec:reblock}). Rows are ordered by F1.}
\end{table}

\begin{table}[h]
\centering
\begin{tabular}{lrrl}
\toprule
Comparison (A vs.\ B) & A right & B right & $p$ \\
\midrule
Decisive/gap vs.\ no selection            & 527 & 150 & $5.2\times10^{-50}$ \\
Decisive/gap vs.\ supervised (200 labels) & 112 &  29 & $1.0\times10^{-12}$ \\
Decisive/gap vs.\ decisive/cumulative 0.8 & 375 & 112 & $3.5\times10^{-34}$ \\
Supervised (200) vs.\ decisive/cumulative 0.8 & 389 & 209 & $1.6\times10^{-13}$ \\
Decisive/cum.\ 0.8 vs.\ no selection      & 468 & 350 & $4.2\times10^{-5}$ \\
Decisive/cum.\ 0.8 vs.\ ordinal/cum.\ 0.8 & 425 & 258 & $1.7\times10^{-10}$ \\
No selection vs.\ ordinal/cum.\ 0.8       & 135 &  86 & $1.2\times10^{-3}$ \\
Ordinal/cum.\ 0.8 vs.\ supervised (10)    & 497 & 453 & $0.16$ \\
\bottomrule
\end{tabular}
\caption{McNemar tests over the pairs on which the two methods disagree, all on the identical
13{,}080-pair candidate set. Counts are pairs the named method got right and the other got
wrong.}
\end{table}

Three readings follow.

\textbf{Recall is effectively constant; every difference is in precision.} Across the table
recall spans 0.988--0.996 while precision spans 0.750--0.896. Attribute selection does not
help the matcher find matches it was missing --- it stops the matcher from asserting matches
that are not there. The one systematic cost is at the aggressive end: the configurations
showing barely more than one attribute give up 18--19 of 2{,}224 matches relative to showing
everything.

\textbf{Elicitation decides whether adaptation helps at all.} At the same selection rule, the
decisive formulation beats no selection ($p = 4.2\times10^{-5}$) while the ordinal
formulation is significantly \emph{worse} than no selection ($p = 1.2\times10^{-3}$). The
pipeline, sampler, transfer model, rule, and matcher are identical between them; only the
profiling question differs. The ordinal formulation also barely selects, retaining all four
attributes on 76.8\% of pairs and averaging 3.77 of 4.

\textbf{The attribute the method learns to discard is \texttt{venue}.} The reason is visible
in the data: the two sources render the same venue very differently ---
\texttt{"SIGMOD Record"} against
\texttt{"The VLDB Journal --- The International Journal on\ldots"} --- so venue disagreement
is evidence of formatting, not of non-identity. The ordinal formulation scores \texttt{venue}
in isolation, where it looks moderately informative, and keeps it.

\subsection{The Selection Rule Matters as Much as the Elicitation}
\label{sec:policy-results}

The cumulative rule's 0.8 target has been in place since the pipeline was written and had
never been examined. Sweeping it downward, with the floor at one attribute and everything
else held fixed, changes the outcome by more than any methodological variable in this report.

\begin{table}[h]
\centering
\begin{tabular}{lrrrrr}
\toprule
Rule & Attrs/pair & FP & FN & F1 & \texttt{title} only \\
\midrule
cumulative 0.8 (default) & 2.24 & 535 & 10 & 0.8904 &  0.0\% \\
cumulative 0.6           & 1.66 & 434 & 13 & 0.9082 & 34.5\% \\
cumulative 0.5           & 1.27 & 293 & 20 & 0.9337 & 73.2\% \\
cumulative 0.4           & 1.18 & 275 & 26 & 0.9359 & 82.7\% \\
gap (no threshold)       & 1.42 & \textbf{255} & 27 & \textbf{0.9397} & 74.9\% \\
\midrule
supervised, fixed \texttt{[title]} & 1.00 & 339 & 26 & 0.9233 & 100.0\% \\
\bottomrule
\end{tabular}
\caption{DBLP--ACM, decisive profiling throughout, varying only how the ranked importance
vector is turned into a subset. The last row is the labelled global baseline, shown as the
limiting case of one attribute on every pair.}
\end{table}

At the shipped default of 0.8, adaptive selection \emph{loses} to the labelled global
baseline ($p = 1.6\times10^{-13}$). At 0.4 it wins ($66/2$, $p = 1.6\times10^{-17}$). The
comparison between adaptive and global selection therefore turned on an untuned
hyperparameter rather than on the methods.

Two further observations. The best configuration is at neither extreme: showing everything
(4.00 attributes) and showing one fixed attribute (1.00) are both beaten by rules that show
roughly 1.2--1.4. And the cumulative targets 0.4 and 0.5 are statistically indistinguishable
from each other ($55/67$, $p = 0.32$), so the effect is a plateau rather than a single lucky
value, while the steps above them are individually significant.

The \texttt{gap} rule is the practically important row. It has no threshold, so nothing about
it was chosen by looking at these results, and it reaches the top of the range anyway. Its
advantage over the best hand-tuned target is thin ($50/31$, $p = 0.045$) and should not be
claimed as superiority; the defensible claim is that it matches a tuned threshold without
requiring one. Since it also does not saturate on wide schemas
(Section~\ref{sec:policy}), it is the rule used for the Amazon--Walmart runs now in progress.

\subsection{Why Adaptation Beats a Fixed Subset}
\label{sec:mechanism}

The labelled global baseline reduces every record to \texttt{title}, which is also what the
adaptive method chooses for three quarters of pairs. Comparing the two therefore isolates
adaptation itself: the same attribute vocabulary, applied uniformly against
pair-by-pair.

\begin{table}[h]
\centering
\begin{tabular}{lrrr}
\toprule
Pair group & Pairs & FP, adaptive & FP, fixed \texttt{[title]} \\
\midrule
Adaptive also chose \texttt{[title]} & 9{,}803 & 118 & 118 \\
Adaptive chose more                  & 3{,}277 & \textbf{137} & 221 \\
\midrule
\quad \texttt{[title, authors, year]} & 2{,}257 & 101 & 146 \\
\quad \texttt{[title, year]}          &    841 &  \textbf{13} &  58 \\
\quad \texttt{[title, authors]}       &    179 &  23 &  17 \\
\bottomrule
\end{tabular}
\caption{DBLP--ACM, all 13{,}080 candidate pairs, comparing Step 1.5 decisive under the
\texttt{gap} rule against the 200-label global baseline over the identical candidate set.}
\end{table}

Where the two select the same attribute they produce identical results --- 118 false
positives each --- which is what the shared matcher cache guarantees and which confirms the
comparison is clean. The entire gain comes from the 3{,}277 pairs where the adaptive method
kept more, cutting false positives from 221 to 137.

The clearest single effect is the 841 pairs given \texttt{[title, year]}: false positives
fall from 58 to 13. These are pairs whose titles are similar enough to mislead a
title-only matcher, and where the year settles the question. A global method cannot express
this, because adding \texttt{year} to every pair is what the cumulative-0.8 configuration
effectively does, and that costs more elsewhere than it saves here.

\subsection{Amazon--Walmart: The Baseline}
\label{sec:amazon}

Amazon--Walmart is the widest schema available and the only one of the three whose
no-selection baseline is far from saturated. Its baseline has now been measured; the three
selection configurations were still running at the time of writing.

\begin{table}[h]
\centering
\begin{tabular}{lrrrrrrr}
\toprule
Method & Attrs/pair & TP & FP & FN & F1 & Precision & Recall \\
\midrule
No selection (\texttt{full}) & 12.00 & 1010 & 1018 & 35 & 0.6573 & 0.4980 & 0.9665 \\
Step 1.5, decisive (\texttt{gap}) & --- & --- & --- & --- & --- & --- & --- \\
Step 1, supervised & --- & --- & --- & --- & --- & --- & --- \\
Step 1.5, ordinal (\texttt{gap}) & --- & --- & --- & --- & --- & --- & --- \\
\bottomrule
\end{tabular}
\caption{Amazon--Walmart, 88{,}296 candidate pairs, 1{,}143 ground-truth matches, blocking
completeness 0.9143 so recall is bounded above by that figure. Runs in progress.}
\end{table}

Two things are already established by this row.

\textbf{This is the benchmark with headroom.} Baseline precision across the three datasets
is 0.4980 here, against 0.7729 on DBLP--ACM and 0.9469 on Fodors--Zagat: close to half the
matches the matcher asserts are wrong. Recall is 0.9665 against a blocking ceiling of
0.9143, so the matcher recovers nearly every candidate available to it and the entire
deficit is in precision --- the quantity attribute selection was shown to improve in
Section~\ref{sec:results-dblp}. It is also the dataset where condensation is worth most in
absolute terms: 1{,}035 prompt tokens per pair at the full representation, against 306 on
DBLP--ACM.

\textbf{Dropping a noise column improved blocking, not just the prompt.} The
\texttt{original\_id} field is a per-table row counter running $1\ldots22{,}074$ and
$1\ldots2{,}554$, sharing no key across the tables. Excluding it raised pair completeness
from 0.9099 to 0.9143, recovering five ground-truth pairs. Because each column is encoded
independently and the column vectors are concatenated, that field occupied a full block of
every record embedding and displaced real matches from the candidate set. No amount of
attribute selection could have recovered them, since Step~1.5 does not re-block --- the
repair had to happen before embedding.

\section{Discussion}
\label{sec:discussion}

\textbf{Pair-adaptive selection outperforms both no selection and a labelled global
baseline.} On DBLP--ACM it improves F1 from 0.8706 to 0.9397 over showing the matcher
everything ($p = 5.2\times10^{-50}$) while showing it 65\% fewer attributes, and beats a
baseline trained on 400 labelled pairs over the identical candidate set
($p = 1.0\times10^{-12}$) while using none.

\textbf{How importance is elicited decides whether adaptation helps.} The ordinal
formulation, on the same pipeline and rule, is significantly worse than no selection at all.
Asking the model to score each attribute in isolation produces a nearly flat importance
vector that no selection rule can act on; asking it to name a sufficient subset forces the
comparison between attributes that the task actually requires. This is the most transferable
finding here, and it is a property of the elicitation, not of the dataset.

\textbf{The rule that consumes the importance vector is as consequential as the vector.} The
same profiling output, read by the default rule, loses to the global baseline; read by
\texttt{gap}, it beats it. Any claim about pair-adaptive selection that does not state its
selection rule is under-specified. That the pipeline's default was both untuned and, on wide
schemas, structurally incapable of condensing at all is the most actionable finding in this
report.

\textbf{The benefit is conditional on the benchmark having headroom.} On Fodors--Zagat, where
the no-selection baseline already reaches F1 0.9596, no method differs from any other.
DBLP--ACM has headroom, and it rewards adaptation even though \texttt{title} nearly
identifies a paper on its own. That the margin over a fixed subset is nonetheless modest ---
141 discordant pairs in 13{,}080 --- is what a near-unique key predicts. A wide schema with
no such key is the sharper test.

\section{Limitations}

\textbf{Amazon--Walmart is incomplete.} Its no-selection baseline is measured
(Section~\ref{sec:amazon}) and confirms the dataset has the headroom the other two lack,
but the three selection configurations were still running when this report was written. The
central claim of this report --- that pair-adaptive selection beats both no selection and a
labelled global baseline --- therefore rests on DBLP--ACM alone, a dataset whose
\texttt{title} field is close to a unique key. This remains the main gap in the evidence.

\textbf{The selection rule was examined on the evaluation set.} The sweep in
Section~\ref{sec:policy-results} was read on the same pairs it reports over, with no held-out
split. It identifies a band worth selecting from, not a defensible value, and no single
cumulative target from it should be quoted. This is why \texttt{gap} is preferred: having no
threshold, it is not exposed to this objection, though it was still chosen after seeing these
results rather than before.

\textbf{Single matcher draw.} Each configuration was run once. The paired design makes the
comparisons valid within that draw --- differing pairs are counted directly, agreeing pairs
share one answer --- but it does not establish that the effect sizes replicate exactly under
a different draw. The large DBLP--ACM effects, with hundreds of discordant pairs, are
unlikely to reverse; the $\texttt{gap}$-versus-tuned-threshold comparison, with 81 discordant
pairs and $p = 0.045$, could.

\textbf{The supervised baseline differs in supervision, not only in method.} It uses 400
labelled pairs; the adaptive configurations use none. Comparing them measures method and
label budget together. Its selected subset also moves with the label budget --- 
\texttt{[title, authors]} at 10 and 50 labels, \texttt{[title]} at 200 --- so no single
sample size characterises it.

\textbf{The feature vector is weaker than it appears.} Of the six per-attribute features the
transfer model consumes, \texttt{semantic\_similarity} is disabled and constant at zero, and
\texttt{is\_numeric} is constant per column on both datasets reported here. On the widest
schema \texttt{exact\_match} fires on 0.06\% of cells. Effectively three features carry
signal. The vector also has no notion of value rarity, which is the standard discriminative
signal in entity resolution: matching on a rare title token and on a common one are
represented identically.

\textbf{One sampler.} All Step~1.5 runs use the \texttt{diversity} profile sampler. No claim
about sampler choice is made here.

\section{Remaining Work}

\begin{enumerate}
\item Complete the three outstanding Amazon--Walmart configurations. With the baseline now
      at F1 0.6573 and precision 0.4980, this is the decisive test of whether the result
      generalizes to a wide schema with no near-unique key.
\item Select the selection rule and any threshold on held-out pairs rather than on the
      evaluation set. \texttt{gap} removes the threshold but not the rule choice.
\item Repeat the DBLP--ACM comparison under an independent matcher draw, to confirm the
      effect sizes.
\item Add a value-rarity feature and re-measure; the current feature set gives the transfer
      model no way to distinguish a distinctive match from a generic one.
\item Re-establish the sampler comparison on the corrected pipeline, on DBLP--ACM rather than
      Fodors--Zagat, where the effect can be resolved.
\end{enumerate}

\section{Summary}

On DBLP--ACM, pair-adaptive attribute selection driven by a subset-elicitation prompt and the
\texttt{gap} selection rule reaches F1 0.9397 using 1.42 of 4 attributes per pair. It beats
showing the matcher every attribute (0.8706, $p = 5.2\times10^{-50}$) and beats a global
baseline trained on 400 labelled pairs that reduces every record to \texttt{title} (0.9233,
$p = 1.0\times10^{-12}$), while using no labels itself. All gains are in precision; recall
falls slightly at the aggressive end. On Fodors--Zagat, where the no-selection baseline
already reaches F1 0.9596, no method is distinguishable from any other, and the defensible
claim is a 46\% reduction in attributes at unchanged accuracy.

Two conditions govern the result. The LLM must be asked for a decision rather than for a
score: the same pipeline driven by an ordinal scoring prompt is significantly worse than no
selection at all. And the rule turning importance into a subset must be able to express small
subsets: under the pipeline's own default the same method loses to the global baseline it
otherwise beats. Neither condition is visible from the aggregate metric alone, which is why
both are reported here.

\end{document}
